\documentclass[12pt, a4paper]{article} % Clase de documento: 'article' es estándar para trabajos cortos. 12pt es el tamaño de letra. A4 es el formato de papel.

% --- PAQUETES ESENCIALES ---
\usepackage[utf8]{inputenc} % Para la codificación de caracteres, esencial para acentos y eñes
\usepackage[T1]{fontenc}    % Fuente moderna para PDFs
\usepackage[spanish]{babel} % Soporte para idioma español (guionización, fechas, etc.)
\usepackage{amsmath}        % Para matemáticas avanzadas (si fuera necesario)
\usepackage{graphicx}       % Para insertar imágenes
\usepackage{booktabs}       % Mejora el aspecto de las tablas (líneas más profesionales)
\usepackage{geometry}       % Para configurar los márgenes
\usepackage{hyperref}       % Para crear enlaces en el índice y referencias (¡muy recomendado!)
\usepackage{fancyhdr}       % Para personalizar encabezados y pies de página (opcional)
\usepackage[spanish]{babel}
% ¡Asegúrate de que xcolor esté aquí!
\usepackage[table]{xcolor} % 'table' es útil para compatibilidad con listings/tablas
\usepackage{listings}       % Paquete para el formato de código fuente
% --- MODIFICACIÓN DE LA ETIQUETA 'LISTING' A 'LISTADO' ---
\addto{\captionsspanish}{\renewcommand{\lstlistingname}{Listado}}
\addto{\captionsspanish}{\renewcommand{\lstlistlistingname}{Listado de Códigos}}
% --- DEFINICIÓN DE COLOR GRIS Y ESTILO DE CÓDIGO ---

% Definimos un color gris claro para el fondo (usando el modelo RGB)
% Los valores entre 0 y 1. (0.95, 0.95, 0.95 es un gris muy claro)
\definecolor{GrisClaro}{rgb}{0.95, 0.95, 0.95}

% Configuramos el estilo por defecto para todos los bloques de código (lstlisting)
\lstset{
    language=TeX, % Lenguaje por defecto (puedes cambiarlo, ej: Python, C, Java)
    backgroundcolor=\color{GrisClaro}, % Fondo gris claro
    basicstyle=\ttfamily\small,      % Tipo de fuente: typewriter (\ttfamily) y pequeño (\small)
    breaklines=true,                 % Permite saltos de línea largos
    showstringspaces=false,          % No mostrar espacios como caracteres
    numbers=left,                    % Numeración de líneas a la izquierda
    numberstyle=\tiny\color{gray},   % Estilo del número de línea
    frame=single,                    % Dibuja un marco alrededor del código
    framesep=5pt,                    % Espacio entre el marco y el código
    rulesepcolor=\color{black},      % Color de las líneas del marco
    keywordstyle=\color{blue}\bfseries, % Palabras clave en azul y negrita
    commentstyle=\color{green!60!black}\itshape, % Comentarios en verde oscuro y cursiva
    stringstyle=\color{orange},      % Cadenas de texto en naranja
}

% --- LENGUAJES PERSONALIZADOS PARA listings (Dart, JavaScript) ---
\lstdefinelanguage{Dart}{
    keywords={import, class, abstract, static, final, const, void, return, if, else, async, await, Future, bool, String, int, double, true, false, null, this, super, extends, new, try, catch, on, in, is, as, override, enum, switch, case, default, late, factory, required, rethrow},
    keywordstyle=\color{blue}\bfseries,
    ndkeywords={SecurityContext, HttpClient, IOClient, X509Certificate, HandshakeException, ApiClient, CertificatePinning, Uint8List, debugPrint, utf8},
    ndkeywordstyle=\color{purple}\bfseries,
    sensitive=true,
    comment=[l]{//},
    morecomment=[s]{/*}{*/},
    morestring=[b]",
    morestring=[b]',
}

\lstdefinelanguage{JavaScript}{
    keywords={var, let, const, function, return, if, else, try, catch, new, this, typeof, null, true, false},
    keywordstyle=\color{blue}\bfseries,
    ndkeywords={Java, Interceptor, Module, console, send, implementation, attach, findExportByName},
    ndkeywordstyle=\color{purple}\bfseries,
    sensitive=true,
    comment=[l]{//},
    morecomment=[s]{/*}{*/},
    morestring=[b]",
    morestring=[b]',
}

% --- CONFIGURACIÓN DE MARGENES ---
\geometry{
 a4paper,
 total={170mm,257mm},
 left=25mm,
 right=25mm,
 top=30mm,
 bottom=30mm,
 }

% --- CONFIGURACIÓN DE ENCABEZADOS Y PIES (OPCIONAL) ---
\pagestyle{fancy}
\fancyhf{} % Limpiar encabezados y pies
\renewcommand{\headrulewidth}{0pt} % Eliminar línea en el encabezado
\rfoot{\thepage} % Número de página a la derecha en el pie
\lfoot{ Seguridad de la Información / Miguel Ángel Gutiérrez Gómez} % Texto en el pie de página izquierdo

% --- DATOS DEL DOCUMENTO (PARA LA PORTADA) ---
\newcommand{\nombreuniversidad}{Universidad Politécnica de Chiapas}
\newcommand{\nombrecarrera}{Ingeniería en Software}
\newcommand{\nombremateria}{Seguridad de la Información}
\newcommand{\nombreprofesor}{MC. José Alonso Macías Montoya}
\newcommand{\corte}{Corte 1}
\newcommand{\actividad}{Práctica - Móvil 4 \\ Certificate Pinning (Anclaje de Certificado) \\ Mitigación de ataques Man-in-the-Middle}
\newcommand{\nombres}{Miguel Ángel Gutiérrez Gómez}
\newcommand{\matricula}{233382}
\newcommand{\fechaentrega}{\today} % \today inserta la fecha actual automáticamente
\newcommand{\ciudad}{Suchiapa, Chiapas}

\usepackage[utf8]{inputenc} % ¡ESTA LÍNEA ES FUNDAMENTAL!
\usepackage[T1]{fontenc}
\usepackage[spanish]{babel}
% --- INICIO DEL DOCUMENTO ---
\begin{document}

% --- PORTADA ---
\begin{titlepage}
    \centering
    \vspace*{1cm} % Espacio vertical
    {\Huge\bfseries \nombreuniversidad \par} % Título principal: Nombre de la Universidad
    \vspace{0.5cm}
    {\Large \nombrecarrera \par} % Nombre de la Carrera

    \vspace{2cm}
    \rule{\textwidth}{1.5pt} % Línea gruesa
    \vspace{0.4cm}
    {\Large\bfseries \actividad \par} % Título del Trabajo
    \vspace{0.4cm}
    \rule{\textwidth}{1.5pt} % Línea gruesa

    \vspace{3cm}

    % --- DATOS DEL ALUMNO Y CURSO (EN UNA SOLA COLUMNA) ---
    % Se eliminaron los entornos minipage y \hfill
    \begin{flushleft}
        \centering % Centra las líneas dentro del entorno flushleft
        \textbf{Materia:} \nombremateria \\
        \textbf{Profesor:} \nombreprofesor \\
        \textbf{Corte:} \corte
    \end{flushleft}

    \vspace{0.5cm} % Pequeño espacio de separación

    \begin{flushleft}
        \centering % Centra las líneas dentro del entorno flushleft
        \textbf{Alumno:} \nombres \\
        \textbf{Matrícula:} \matricula \\
    \end{flushleft}
    % --- FIN DE DATOS EN UNA SOLA COLUMNA ---

    \vfill % Empuja el contenido restante al final de la página

    % Fecha y Ciudad
    \vspace{1cm}
    \textbf{\ciudad, a \fechaentrega}

\end{titlepage}

\clearpage % Salto de página para empezar el contenido

% --- ÍNDICE/TABLA DE CONTENIDO ---
\tableofcontents
\clearpage

% --- CUERPO DEL TRABAJO ---

\section{Introducción}
\label{sec:introduccion}
El protocolo HTTPS garantiza que el tráfico entre una aplicación móvil y su
servidor viaje cifrado, pero por sí solo \textbf{no garantiza con quién se
está hablando realmente}. La validación TLS estándar acepta como legítimo
cualquier certificado firmado por alguna de las más de 150 Autoridades
Certificadoras (CA) que el sistema operativo considera de confianza. Basta con
que un atacante logre instalar su propia CA en el dispositivo —algo trivial en
un teléfono que el usuario controla, y que hacen por diseño herramientas como
\textbf{mitmproxy}, \textbf{Burp Suite} o \textbf{Charles Proxy}— para que
pueda descifrar, leer y modificar todo el tráfico sin que la app lo note. Es el
clásico ataque \textit{Man-in-the-Middle} (MITM).

El \textbf{Certificate Pinning} (anclaje de certificado) mitiga este riesgo
invirtiendo la lógica de confianza: en lugar de aceptar a cualquier CA del
sistema, la aplicación lleva embebido el certificado (o su clave pública) del
servidor con el que espera hablar, y \textbf{rechaza la conexión si el
certificado presentado no coincide}.

El objetivo de esta práctica es integrar certificate pinning en la aplicación
\textbf{VisionPrice} (Flutter), que consume una API REST desplegada en Railway
sobre HTTPS, y demostrar empíricamente que la app deja de ser interceptable.

\section{Recursos Utilizados}
\label{sec:recursos}
En esta sección se listan las herramientas y servicios fundamentales para la
realización de la actividad.

\begin{itemize}
    \item \textbf{Cliente móvil:} Flutter 3.44 / Dart 3.12, paquete
    \texttt{http} y la clase \texttt{SecurityContext} de \texttt{dart:io}.
    \item \textbf{Backend:} microservicio \texttt{auth-service}
    (Node.js + Fastify + Prisma) desplegado en \textbf{Railway} con HTTPS y
    certificado emitido por Let's Encrypt.
    \item \textbf{Análisis de certificados:} \texttt{openssl s\_client} para
    extraer la cadena de certificación y calcular las huellas SPKI.
    \item \textbf{Interceptación (verificación):} \textbf{mitmproxy} como
    proxy MITM de prueba.
    \item \textbf{Otros:} Android Studio, VS Code, Git, Overleaf y compilador
    \LaTeX{}.
\end{itemize}

\section{Desarrollo}
\label{sec:desarrollo}
Esta es la parte central del trabajo. Se presentan el análisis previo, la
justificación de la estrategia elegida, la implementación y su verificación.

\subsection{Análisis de la cadena de certificación}
\label{sec:cadena}
El primer paso fue conocer exactamente qué certificado presenta el servidor.
Se consultó con \texttt{openssl}:

\begin{lstlisting}[language=bash, caption={Extracción de la cadena de certificación}]
openssl s_client -connect vsauth-production.up.railway.app:443 \
    -servername vsauth-production.up.railway.app -showcerts
\end{lstlisting}

La cadena obtenida consta de cuatro certificados. Para cada uno se calculó la
huella SHA-256 de su clave pública (\textbf{SPKI}), que es el identificador
estándar recomendado para anclar:

\begin{table}[h]
    \centering
    \caption{Cadena de certificación y huellas SPKI obtenidas}
    \label{tab:cadena}
    \begin{tabular}{clp{5.2cm}}
        \toprule
        \textbf{Nivel} & \textbf{Sujeto} & \textbf{SPKI (SHA-256, base64)} \\
        \midrule
        0 (hoja)  & \texttt{*.up.railway.app}   & \texttt{\scriptsize hORJtOUv6S5sVpxykVdj7ceRoUfeLhICfOGA0n68co0=} \\
        1 (inter.)& \texttt{Let's Encrypt YE1}  & \texttt{\scriptsize brzvtCELCIZUo4sD/qPX0ccRtPsd3DY6RfmxpOU9oB4=} \\
        2 (raíz)  & \texttt{ISRG Root YE}       & \texttt{\scriptsize sCkq5UWXjg+7mKu9lMhhYF5bGLsy7VI/UNW3tccdR7w=} \\
        3 (raíz)  & \texttt{ISRG Root X2}       & \texttt{\scriptsize diGVwiVYbubAI3RW4hB9xU8e/CH2GnkuvVFZE8zmgzI=} \\
        \bottomrule
    \end{tabular}
    \vspace{0.5em}
    \footnotesize Nota: datos reales obtenidos del servidor de producción.
\end{table}

\subsection{Justificación: qué nivel de la cadena anclar}
\label{sec:justificacion}
Ésta es la decisión de diseño más importante de la práctica, y la que
determina si la aplicación seguirá funcionando dentro de tres meses.

El certificado hoja es un \textbf{comodín} (\texttt{*.up.railway.app}) que
\textbf{pertenece a Railway, no al desarrollador}. Let's Encrypt emite
certificados con una validez máxima de 90 días, y Railway los rota sin previo
aviso. Se evaluaron tres alternativas:

\begin{enumerate}
    \item \textbf{Anclar el certificado hoja.} Máxima seguridad, pero la
    aplicación \textbf{dejaría de funcionar por sí sola} en cuanto Railway
    renovara el certificado. Sería necesario publicar una nueva versión en la
    tienda cada pocas semanas: inviable.
    \item \textbf{Anclar el intermedio (Let's Encrypt YE1).} Sobrevive a la
    rotación del certificado hoja. Su validez alcanza hasta septiembre de 2028.
    \item \textbf{Anclar la raíz (ISRG Root X2).} La opción más duradera
    (válida hasta 2032) pero también la más permisiva.
\end{enumerate}

\textbf{Se optó por una estrategia de anclas redundantes}: se embebieron
simultáneamente el \textbf{intermedio} y la \textbf{raíz}. Así, si Let's
Encrypt rotara el intermedio antes de 2028, el ancla de la raíz mantendría la
aplicación operativa. Esta técnica de \textit{pines de respaldo} es la práctica
estándar de la industria.

\textbf{Limitación asumida conscientemente:} un atacante que consiguiera que
Let's Encrypt le emitiera un certificado válido \textbf{para ese mismo dominio}
superaría el filtro. Sin embargo, eso exige comprometer el control del dominio,
lo que representa un ataque de un orden de magnitud muy superior al de
simplemente instalar una CA en el teléfono.

\subsection{Arquitectura de la solución}
\label{sec:arquitectura}
La aplicación sigue \textbf{Clean Architecture}, y su clase \texttt{ApiClient}
ya recibía el cliente HTTP por \textbf{inyección de dependencias}:

\begin{lstlisting}[language=Dart, caption={ApiClient: el cliente HTTP era ya inyectable}]
ApiClient({
  http.Client? client,
  required String baseUrl,
}) : _client = client ?? http.Client(), ...
\end{lstlisting}

Gracias a ese diseño previo, activar el anclaje se resolvió \textbf{creando un
único archivo nuevo} (\texttt{core/network/certificate\_pinning.dart}) y
modificando \textbf{tres líneas} del \textit{composition root}. No se tocó
ningún \textit{data source}, repositorio, caso de uso ni ViewModel: el cambio
quedó totalmente contenido en la capa de infraestructura de red.

\subsection{Implementación}
\label{sec:implementacion}
El mecanismo central es construir un \texttt{SecurityContext} con
\texttt{withTrustedRoots: false}, lo que \textbf{descarta por completo el
almacén de CAs del sistema operativo}, e inyectarle únicamente los
certificados aprobados.

\begin{lstlisting}[language=Dart, caption={Construcción del cliente HTTP anclado}]
static http.Client createClient() {
  if (!kEnabled) return http.Client();

  // withTrustedRoots: false descarta el almacen de CAs del sistema.
  final context = SecurityContext(withTrustedRoots: false);
  context.setTrustedCertificatesBytes(
    utf8.encode('$_letsEncryptYe1\n$_isrgRootX2'),
  );

  final ioClient = HttpClient(context: context)
    // NUNCA devolver true aqui: eso anularia por completo el anclaje.
    ..badCertificateCallback = (X509Certificate cert, String host, int port) {
      debugPrint('Pinning: certificado RECHAZADO para $host:$port');
      return false;
    }
    ..connectionTimeout = const Duration(seconds: 15);

  return IOClient(ioClient);
}
\end{lstlisting}

Y su conexión en el \textit{composition root}:

\begin{lstlisting}[language=Dart, caption={Inyección del cliente anclado en ambos ApiClient}]
final pinnedClient = CertificatePinning.createClient();
_authApiClient = ApiClient(client: pinnedClient, baseUrl: ApiConfig.baseUrl);
_projectsApiClient =
    ApiClient(client: pinnedClient, baseUrl: ApiConfig.projectsBaseUrl);
\end{lstlisting}

\subsubsection{El error más común al implementar pinning}
Numerosos tutoriales implementan el anclaje devolviendo \texttt{true} desde
\texttt{badCertificateCallback}. Eso \textbf{elimina toda la seguridad}: ese
callback solo se invoca cuando la validación \textit{ya falló}, de modo que
devolver \texttt{true} equivale a aceptar cualquier certificado, incluido el
del atacante. En esta implementación devuelve siempre \texttt{false} y se
utiliza únicamente para registrar el rechazo.

\subsection{Verificación empírica}
\label{sec:verificacion}
Para comprobar que el anclaje funciona se ejecutó un programa Dart que realiza
la misma petición con dos clientes distintos: uno anclado y uno normal.

\begin{table}[h]
    \centering
    \caption{Resultados de la verificación del anclaje}
    \label{tab:resultados}
    \begin{tabular}{lcc}
        \toprule
        \textbf{Cliente} & \textbf{Servidor propio (Railway)} & \textbf{Dominio con otra CA} \\
        \midrule
        Con pinning        & Permitido (200 OK) & \textbf{Bloqueado} \\
        Sin pinning (control) & Permitido (200 OK) & Permitido \\
        \bottomrule
    \end{tabular}
    \vspace{0.5em}
    \footnotesize Nota: el bloqueo se produce con \texttt{HandshakeException}.
\end{table}

El detalle relevante es \textbf{en qué momento} falla la conexión: la excepción
es \texttt{HandshakeException}, es decir, la conexión se aborta durante la
negociación TLS, \textbf{antes de transmitir un solo byte de datos de
aplicación}. Ninguna credencial ni token llega jamás al interceptor.

\subsection{Capturas de funcionamiento}

\begin{figure}[h]
    \centering
    \fbox{\parbox[c][5.5cm][c]{0.62\textwidth}{\centering \textit{mitmproxy interceptando el tráfico con el anclaje DESACTIVADO: se observan en claro las peticiones de login y de proyectos.}}}
    \caption{Sin pinning: el tráfico es legible por el proxy.}
    \label{fig:sinpinning}
\end{figure}

\begin{figure}[h]
    \centering
    \fbox{\parbox[c][5.5cm][c]{0.62\textwidth}{\centering \textit{Con el anclaje ACTIVADO: mitmproxy registra el handshake abortado por el cliente y la app no logra conectar.}}}
    \caption{Con pinning: el handshake TLS es rechazado.}
    \label{fig:conpinning}
\end{figure}

\section{Resultados y Conclusiones}
\label{sec:resultados}
En esta sección se presentan los \textbf{resultados clave} derivados del
desarrollo, seguidos por las \textbf{conclusiones}.

\begin{itemize}
    \item \textbf{Resultado 1:} Se integró certificate pinning en VisionPrice
    verificando la cadena contra dos anclas embebidas, sin depender del
    almacén de CAs del sistema operativo.
    \item \textbf{Resultado 2:} Se comprobó empíricamente que un dominio
    servido por una CA distinta es rechazado con \texttt{HandshakeException},
    mientras que el servidor legítimo sigue respondiendo con normalidad.
    \item \textbf{Resultado 3:} El cambio quedó contenido en un archivo nuevo
    y tres líneas del \textit{composition root}, evidenciando el beneficio
    práctico de haber aplicado inyección de dependencias en la capa de red.
    \item \textbf{Conclusión 1:} HTTPS por sí solo protege la
    \textit{confidencialidad} del canal, pero no la \textit{identidad} del
    interlocutor cuando el atacante controla el dispositivo. El pinning es la
    capa que cierra esa brecha.
    \item \textbf{Conclusión 2:} La seguridad y la mantenibilidad están en
    tensión directa. Anclar el certificado hoja es lo más seguro y lo menos
    sostenible; elegir el nivel adecuado de la cadena es una decisión de
    ingeniería, no una receta única.
\end{itemize}

\clearpage % Salto de página para empezar los anexos

\section{Anexos}
\label{sec:anexos}
Los anexos incluyen material complementario que respalda el cuerpo principal
del trabajo.

\subsection{Repositorio del proyecto}
El código fuente completo (aplicación Flutter y microservicios) se encuentra en:
\begin{center}
\url{https://github.com/MiguelAngelGutierrezGomez/VisionPrice}
\end{center}

\subsection{Análisis Crítico: ¿es infalible frente a Frida?}
\label{sec:analisis}
\textbf{No.} El certificate pinning es una defensa \textbf{del lado del
cliente}, y todo lo que se ejecuta en un dispositivo bajo control del atacante
es, en última instancia, modificable. Un atacante con \textbf{Frida} opera
\textit{dentro} del proceso de la aplicación y dispone de varias vías:

\begin{itemize}
    \item \textbf{Hookear el callback de validación.} Interceptar
    \texttt{badCertificateCallback} y forzar que devuelva \texttt{true}. Son
    pocas líneas de JavaScript.
    \item \textbf{Neutralizar la verificación en BoringSSL.} Flutter no usa
    la pila TLS de Android, sino BoringSSL compilada dentro de
    \texttt{libflutter.so}. Existen scripts públicos que localizan por patrón
    la función \texttt{ssl\_crypto\_x509\_session\_verify\_cert\_chain} y la
    parchean en memoria para que siempre apruebe.
    \item \textbf{Reemplazar los certificados embebidos.} Extraer el APK,
    sustituir los PEM anclados por el de su propia CA y re-empaquetar.
\end{itemize}

\begin{lstlisting}[language=JavaScript, caption={Concepto de un bypass con Frida (ilustrativo)}]
// Localiza la verificacion de cadena dentro de libflutter.so y la
// reemplaza por una funcion que siempre devuelve "valido".
Interceptor.replace(direccionDeLaFuncion, new NativeCallback(function () {
    return 1; // 1 = certificado aceptado
}, 'int', ['pointer', 'pointer', 'pointer']));
\end{lstlisting}

\subsubsection{Capas adicionales que se añadirían}
\begin{enumerate}
    \item \textbf{Detección de instrumentación:} búsqueda de
    \texttt{frida-server} (puertos y \textit{named pipes} por defecto), de
    bibliotecas \texttt{frida-gadget} cargadas en memoria y de indicios de
    \textit{root} (binario \texttt{su}, paquetes de Magisk).
    \item \textbf{Anti-debugging nativo (NDK):}
    \texttt{ptrace(PTRACE\_TRACEME)} y lectura de \texttt{TracerPid} en
    \texttt{/proc/self/status} para detectar depuradores adjuntos.
    \item \textbf{Verificación de integridad:} comprobación de la firma del
    APK en tiempo de ejecución para detectar re-empaquetado, y
    \textbf{Play Integrity API} como atestación respaldada por Google.
    \item \textbf{Defensa del lado servidor:} \textbf{ésta es la capa
    decisiva}. Toda regla de negocio crítica debe validarse en el backend,
    junto con detección de anomalías y rotación de tokens. El servidor es el
    único entorno que el atacante no controla.
\end{enumerate}

\subsubsection{Conclusión del análisis}
El valor del pinning no reside en ser inviolable, sino en \textbf{elevar
sustancialmente el costo del ataque}: sin él, interceptar la aplicación
requiere instalar un proxy y una CA, algo que puede hacer cualquier persona
siguiendo un tutorial; con él, se necesita un dispositivo \textit{rooteado},
conocimientos de instrumentación dinámica y de ingeniería inversa de binarios
nativos. Esa diferencia elimina a la inmensa mayoría de los atacantes
oportunistas. La regla permanente es \textbf{no confiar jamás en el
dispositivo} como entorno de ejecución seguro.

\subsection{Código Fuente Adicional}
A continuación, el bloque completo de anclas embebidas y el comando empleado
para regenerarlas cuando la cadena del servidor cambie.

\begin{lstlisting}[language=bash, caption={Comando para recalcular la huella SPKI de un certificado}]
openssl x509 -in cert.pem -pubkey -noout \
  | openssl pkey -pubin -outform der \
  | openssl dgst -sha256 -binary \
  | openssl enc -base64
\end{lstlisting}

\begin{lstlisting}[language=Dart, caption={Interruptor maestro y anclas embebidas (extracto)}]
class CertificatePinning {
  /// Interruptor maestro. En `false` desactiva el anclaje para poder
  /// inspeccionar el trafico con un proxy durante una demostracion.
  static const bool kEnabled = true;

  /// Intermedio: Let's Encrypt YE1 (valido 2025-09-03 -> 2028-09-02).
  static const String _letsEncryptYe1 = '''
-----BEGIN CERTIFICATE-----
MIICizCCAhGgAwIBAgIQXd1w3TH4AchcGGp6BLgK/jAKBggqhkjOPQQDAzAuMQsw
... (contenido completo en el repositorio) ...
-----END CERTIFICATE-----
''';

  /// Raiz: ISRG Root X2 (valida 2026-05-13 -> 2032-09-02). Ancla de
  /// respaldo por si Let's Encrypt rota el intermedio antes de 2028.
  static const String _isrgRootX2 = '''
-----BEGIN CERTIFICATE-----
MIIEcDCCAligAwIBAgIQbI8dxyfHEX97r4U6yYD5zTANBgkqhkiG9w0BAQsFADBP
... (contenido completo en el repositorio) ...
-----END CERTIFICATE-----
''';
}
\end{lstlisting}

\end{document}
